\section{Exact type restrictions and separation of a shared leg}
\label{sec:maps}

We now realize the two families of relations by finite tensor maps.
There are two distinct tasks. A type projection removes off-diagonal
blocks, but leaves the first leg shared. A subsequent separator makes
the surviving blocks independent on all three legs.
Precisely, matched sectors sharing the first leg have the form
\(\sum_h A_h\), with \(A_h\in X\otimes Y_h\otimes Z_h\) and
direct-sum decompositions on the second and third legs. All supported
terms use the same sector label on those two legs, while the space \(X\)
is common to all sectors.

\subsection{The neighboring-length map}

Let \(V_e\) be the homogeneous binary forms of degree \(e\) in \(u,v\),
and let \(W=\operatorname{span}(s,w)\). For degree \(e\ge1\), put
\(s_i=u^{e-i}v^is\) and \(w_i=u^{e-i}v^iw\).
Choose quotient lifts \(s_0,\ldots,s_e,w_e\) and kernel vectors
\(k_i=w_i-s_{i+1}\), \(0\le i<e\).
These form a basis of \(V_e\otimes W\): substitution \((s,w)=(u,v)\)
sends the quotient lifts to distinct monomials and kills the kernel
vectors, which are independent. The kernel is
\((uw-vs)V_{e-1}\).

\begin{lemma}[Exact neighboring-length types]\label{lem:determinant}
For positive integers \(n_+,n_-\) with \(n=n_++n_-\), a specified
restriction of
\((C(r,b)B_X(2))^{\otimes n}\), \(r,b\ge2\), has
\(\binom n{n_+}\) matched sectors on its second and third legs.
Each sector is isomorphic to
\(C(r,b+1)^{\otimes n_+}C(r,b-1)^{\otimes n_-}\).
The first leg is still shared; the sector-specific coordinate permutation
is applied only after separation.
\end{lemma}
\begin{proof}
View \(C(r,b)B_X(2)\) as multiplication
\(V_{r-1}\times(V_{b-1}\otimes W)\to V_{r+b-2}\otimes W\).
Use the preceding bases on its second-input and output spaces, and the
dual output basis on its third leg. Multiplication by
\(u^{r-1-i}v^i\) shifts each monomial index by \(i\).
It preserves the kernel and induces ordinary multiplication on the quotient.
The two diagonal blocks are \(C(r,b+1)\) and \(C(r,b-1)\).
The only off-diagonal terms are
\begin{equation}\label{eq:off-diagonal}
 x_i q_b k^*_{b-1+i},\qquad 0\le i<r-1,
\end{equation}
where \(q_b=w_{b-1}\) is the last input quotient lift and the star
denotes an output-dual coordinate.
Indeed, the image of that lift has index \(b-1+i\). Unless this is the
last output index, \(w_{b-1+i}=s_{b+i}+k_{b-1+i}\).
All other quotient lifts stay in the quotient, and a kernel vector maps
to a kernel vector. This proves \eqref{eq:off-diagonal}.

In \(n\) tensor positions, retain words with exactly \(n_-\) kernel
labels on both the second and third legs. Each off-diagonal term increases
the output kernel count by one, and no term decreases it. Equal retained
counts therefore remove every off-diagonal term. In a surviving term the
two sector labels agree at each position. Choosing the \(n_+\) quotient
positions gives the asserted number of sectors and their tensors.
\end{proof}

This is an exact restriction after a fixed basis change. In particular,
it incurs no interpolation degree.

\subsection{The three-sector map}

\begin{lemma}[Exact three-sector types]\label{lem:tripling}
For \(r\ge2\) and \(3\mid n\), a coordinate restriction of
\(C(r,4r-1)^{\otimes n}\) has
\(n!/(n/3)!^3\) matched sectors on its second and third legs.
Each sector contains \(2n/3\) factors isomorphic to \(C(r,r)\) and
\(n/3\) factors isomorphic to \(C(r,r)\) with its second and third
legs exchanged. The first leg remains shared.
The sector-specific coordinate permutations are applied only after separation.
\end{lemma}
\begin{proof}
A supported term has \(0\le i<r\), \(0\le y<4r-1\), and \(z=i+y\).
Partition the index sets into the following intervals, with endpoints
included:
\[
\begin{array}{c|ccc}
 &\text{left}&\text{middle}&\text{right}\\ \hline
Y &[0,r-1]&[r,3r-2]&[3r-1,4r-2]\\
Z &[0,2r-2]&[2r-1,3r-2]&[3r-1,5r-3].
\end{array}
\]
A left \(Y\)-index forces a left \(Z\)-index, and a right \(Y\)-index
forces a right \(Z\)-index. Every unmatched term therefore has a middle
\(Y\)-label and an outer \(Z\)-label. Retain \(n/3\) labels of each
kind on both legs in a tensor power. Equality of the middle counts
eliminates every unmatched term; the remaining labels agree positionwise.

The outer blocks are translates of \(C(r,r)\).
The middle block has terms
\[
 x_{r-1-i}\,y_{r+i+j}\,z_{2r-1+j},\qquad 0\le i,j<r.
\]
Reversing the first index and exchanging the last two legs identifies it
with \(C(r,r)\). These identifications are statements about each sector
separately. We do not apply different first-leg maps to blocks while that
leg is shared.
\end{proof}

\subsection{A Fourier separator}

The following form includes an auxiliary dimension and a phase slope.
The main construction uses only \(K_0=M\), \(c=1\).

\begin{proposition}[Shared-leg separator]\label{prop:separator}
Let
\(A=\sum_{h=0}^{M-1}A_h\), with
\(A_h\in X\otimes Y_h\otimes Z_h\), where the \(Y_h\) and \(Z_h\)
are pairwise disjoint sectors. For positive integers \(K_0,c\), put
\(L=K_0+c(M-1)\). There is an explicit degeneration
\begin{equation}\label{eq:separator}
 L A\degen\bigoplus_{h=0}^{M-1}(A_h B_X(K_0)).
\end{equation}
Its polynomial has exponent support contained in
\(\{cd^2:0\le d\le d_*\}\), where
\(d_* =\min(M-1,\lfloor(K_0-1)/c\rfloor)\).
In particular, \(K_0=M,c=1\) uses \(2M-1\) source copies and has
degree at most \((M-1)^2\).
Permuting all three legs gives the same statement with any chosen shared leg.
\end{proposition}
\begin{proof}
Choose a primitive \(L\)-th root \(\zeta\). In source copy \(\ell\),
\(0\le\ell<L\), make the substitutions
\begin{align*}
 x_a^{(\ell)}&\longmapsto\sum_{g=0}^{M-1}\zeta^{\ell cg}X_{a,g},\\
 y_{h,b}^{(\ell)}&\longmapsto
       \sum_{u=0}^{K_0-1}\zeta^{\ell(u-ch)}Y_{h,b,u},\\
 z_{h,d}^{(\ell)}&\longmapsto
       L^{-1}\sum_{v=0}^{K_0-1}\zeta^{-\ell v}Z_{h,d,v}.
\end{align*}
Summing the copies retains exactly
\(u-v+c(g-h)\equiv0\pmod L\). Its absolute value is at most \(L-1\),
so the congruence is the equality \(u-v+c(g-h)=0\).
Multiply the new variables by the respective powers of \(t\)
\[
 t^{cg^2},\qquad t^{2hu-ch^2},\qquad t^{-2hv}.
\]
Every surviving term has total exponent
\(cg^2+2h(u-v)-ch^2=c(g-h)^2\).
It is zero exactly when \(g=h\) and \(u=v\). These terms are precisely
the target in \eqref{eq:separator}, whose first variables \(X_{a,h}\)
are now independent for different \(h\).
Moreover \(c|g-h|=|u-v|\le K_0-1\), proving the exponent support.
\end{proof}

Only after this separator has been applied do we reorder tensor positions
or use sector-specific first-leg reversals. Applying the maps on all six
permutations balances the exchanged factors in Lemma~\ref{lem:tripling}:
composition with the transposition \(\sigma\) of its last two legs
permutes \(S_3\). Explicitly, for \(C=C(r,r)\),
\[
 \bigotimes_{\pi\in S_3}
 \bigl((C^\pi)^{\otimes 2n/3}((C^\sigma)^\pi)^{\otimes n/3}\bigr)
 \cong Q_{r,r}^{\otimes n}.
\]
This uses permutations of tensor factors within each global leg; it
does not swap two global legs in only part of a tensor.
Thus every separated balanced sector is exactly the prescribed product
of the \(Q_{r,b}\), up to a specified coordinate permutation.

Further finite refinements of this separator are proved in
Appendix~\ref{subsec:sparse}. They are not used in the main extraction.
